A Kaluza-type lift is a third way of turning forces into geometry. The key idea is different from both Hertz and Eisenhart:

$$ \boxed{\text{the force is encoded in off-diagonal components of a higher-dimensional metric}.} $$

The prototype is Kaluza's five-dimensional treatment of electromagnetism.

1. Start with a charged particle

In four-dimensional spacetime, a particle of mass \(m\) and charge \(q\) has

$$ L = \frac{m}{2}g_{\mu\nu}\dot x^\mu\dot x^\nu + q A_\mu\dot x^\mu . $$

The electromagnetic potential appears linearly in the velocity.

The corresponding equation is the Lorentz-force equation

$$ m\frac{D u^\mu}{d\tau} = q F^\mu{}_{\nu}u^\nu, \qquad F_{\mu\nu} = \partial_\mu A_\nu-\partial_\nu A_\mu . $$

The question is:

Can the \(qA_\mu\dot x^\mu\) term itself arise from ordinary geodesic motion?

The Kaluza answer is yes.

2. Add one dimension

Introduce an extra coordinate \(y\) and consider

$$ X^A=(x^\mu,y). $$

Take the five-dimensional metric

$$ \boxed{ ds_5^2 = g_{\mu\nu}dx^\mu dx^\nu + \phi^2 \left(dy+\kappa A_\mu dx^\mu\right)^2 } $$

where \(\phi\) is a scalar field and \(\kappa\) is a coupling constant.

Expanding,

$$ ds_5^2 = \left(g_{\mu\nu} +\kappa^2\phi^2 A_\mu A_\nu\right)dx^\mu dx^\nu + 2\kappa\phi^2A_\mu dx^\mu dy + \phi^2dy^2 . $$

The crucial term is

$$ \boxed{ g_{\mu y}\sim A_\mu. } $$

So the electromagnetic potential is literally an off-diagonal component of the higher-dimensional metric.

That is the defining Kaluza idea.

3. How the Lorentz force appears

Consider the five-dimensional geodesic Lagrangian

$$ L_5= \frac12G_{AB}\dot X^A\dot X^B. $$

Substitution gives

$$ L_5= \frac12g_{\mu\nu}\dot x^\mu\dot x^\nu + \frac12\phi^2 \left( \dot y+\kappa A_\mu\dot x^\mu \right)^2. $$

Because \(y\) does not appear explicitly,

$$ p_y= \frac{\partial L_5}{\partial\dot y} = \phi^2 \left( \dot y+\kappa A_\mu\dot x^\mu \right) $$

is conserved.

Call this constant \(p_y=k\). Then

$$ \dot y+\kappa A_\mu\dot x^\mu = \frac{k}{\phi^2}. $$

Substituting this back into the \(x^\mu\) equations produces, schematically,

$$ \boxed{ \frac{D u^\mu}{d\tau} \sim k\kappa F^\mu{}_\nu u^\nu +\text{terms involving }\phi . } $$

If \(\phi\) is constant, this is precisely the Lorentz-force equation after identifying

$$ q/m\sim k\kappa . $$

Thus

$$ \boxed{ \text{5D geodesic} \quad\longrightarrow\quad \text{4D charged-particle motion}. } $$

This is conceptually very different from Eisenhart.

4. Compare the three constructions

Now the distinction becomes particularly clean.

Hertz
$$ \boxed{ \text{force} \longrightarrow \text{extra mechanical coordinates} } $$

and the physical dynamics is a projection of free geodesic motion.

Eisenhart
$$ \boxed{ V \longrightarrow g_{tt} } $$

and Newtonian trajectories are null geodesics.

Kaluza
$$ \boxed{ A_\mu \longrightarrow G_{\mu y} } $$

and charged trajectories are projections of higher-dimensional geodesics.

So:

$$ \begin{array}{c|c} \text{Construction}&\text{What carries the dynamics}\\ \hline \text{Hertz}&\text{extra configuration-space structure}\\ \text{Eisenhart}&g_{tt}\\ \text{Kaluza}&g_{\mu y} \end{array} $$
5. Why the gauge transformation becomes geometrical

This is perhaps the most beautiful aspect.

Electromagnetism has the gauge transformation

$$ A_\mu\rightarrow A_\mu+\partial_\mu\lambda. $$

Now make the coordinate transformation

$$ \boxed{ y\rightarrow y-\kappa\lambda(x). } $$

Then

$$ dy\rightarrow dy-\kappa\,d\lambda, $$

so

$$ dy+\kappa A_\mu dx^\mu $$

is unchanged:

$$ dy-\kappa d\lambda + \kappa(A+d\lambda) = dy+\kappa A. $$

Therefore

$$ A_\mu\rightarrow A_\mu+\partial_\mu\lambda $$

is simply the manifestation of a coordinate transformation in the higher-dimensional space.

This is the fundamental Kaluza insight:

$$ \boxed{ \text{gauge symmetry} \sim \text{higher-dimensional coordinate symmetry}. } $$
6. The modern geometric interpretation

The construction is best understood using a principal \(U(1)\) bundle.

The extra coordinate \(y\) is a coordinate along the fiber, and

$$ \boxed{ \theta=dy+\kappa A } $$

is a connection one-form.

Its curvature is

$$ d\theta = \kappa\,dA = \kappa F. $$

Thus

$$ \boxed{ A = \text{connection}, \qquad F = \text{curvature}. } $$

This is exactly the connection/curvature distinction you have been emphasizing in your discussions of gauge structure.

The Kaluza metric is essentially

$$ \boxed{ G = g_{\text{base}} + \phi^2\theta^2. } $$

So the higher-dimensional metric contains simultaneously:

$$ \begin{array}{ccl} g_{\mu\nu}&\rightarrow&\text{gravity},\\ A_\mu&\rightarrow&U(1)\text{ connection},\\ \phi&\rightarrow&\text{scalar field}. \end{array} $$
7. Non-Abelian generalization

This becomes even more interesting.

Replace \(U(1)\) by a Lie group \(G\), with gauge potential

$$ A=A_\mu^a T_a\,dx^\mu. $$

The curvature is

$$ F=dA+A\wedge A, $$

or

$$ F_{\mu\nu}^a = \partial_\mu A_\nu^a -\partial_\nu A_\mu^a + f^a{}_{bc}A_\mu^bA_\nu^c. $$

A higher-dimensional space is constructed as a principal \(G\)-bundle over spacetime.

Then schematically,

$$ \boxed{ \text{higher-dimensional metric} \longrightarrow g_{\mu\nu}+A_\mu^a+\text{scalar fields}. } $$

This is the geometric origin of the modern Kaluza–Klein picture.

8. Kaluza versus Eisenhart: an important distinction

There is a particularly useful conceptual difference.

In Eisenhart,

$$ V(q,t) $$

is put into the metric itself as

$$ g_{tt}\sim -V. $$

Thus the potential is a metric component on spacetime.

In Kaluza,

$$ A_\mu $$

appears in an off-diagonal metric component

$$ G_{\mu y}\sim A_\mu. $$

But the physically important electromagnetic field is

$$ F=dA, $$

which is the curvature of the connection.

Thus:

$$ \boxed{ \begin{array}{rcl} \text{Eisenhart} &: & V\rightarrow\text{metric}\\ \text{Kaluza} &: & A\rightarrow\text{connection}\\ &&F\rightarrow\text{curvature}. \end{array}} $$
9. There is a very nice unified picture

We can now put all three constructions into one framework:

$$ \boxed{ \begin{array}{ccc} \text{Hertz} & \text{Eisenhart} & \text{Kaluza} \\[3pt] \downarrow&&\downarrow \\[-2pt] \text{extra configuration} & \text{extra null direction} & \text{fiber direction} \\ \downarrow& \downarrow& \downarrow \\ \text{geodesics} & \text{null geodesics} & \text{geodesics} \\ \downarrow& \downarrow& \downarrow \\ \text{mechanics} & \text{Newtonian mechanics} & \text{charged mechanics} \end{array}} $$

But there is an even deeper connection:

Eisenhart and Kaluza can actually be combined.

One can construct a Bargmann/Kaluza-type metric in which both a scalar potential and a gauge potential occur:

$$ \boxed{ ds^2 = g_{ij}dq^i dq^j + 2dt\left(ds+A_i\,dq^i-\Phi\,dt\right) } $$

(up to convention and normalization).

Then the null geodesics reproduce a particle subject simultaneously to scalar and magnetic forces.

This is especially relevant to your interests in stochastic mechanics with gauge fields: the object that naturally appears is no longer merely a scalar potential \(V\), but a connection \(A\), and its curvature \(F\) produces the magnetic force. In geometric language, Eisenhart handles the scalar part while the Kaluza/Bargmann structure handles the connection part.

That combined construction is probably the most useful next step if we want to connect these ideas to Nelson mechanics on manifolds, magnetic fields, and stochastic parallel transport.
